\scanpage{54}
\uline{Catégorie modérée}

A) Catégorie $\mathcal M$ avec foncteur
\[
X\longmapsto|X|\qquad
\varphi:\mathcal M\longrightarrow(\text{Espaces compacts})
\]
[M0) $|\;|$ foncteur ``transportable'']

M1) \deleted{$\unsure{\varphi}$} Dans $\mathcal M$ les $\varprojlim$ finies existent, et $\varphi$ est exact à g.

M2) Le foncteur $|\;|$ \deleted{$\unsure{\varphi}$} est \deleted{\unsure{fidèle}}\inserted{conservatif}
[$\Longrightarrow$ le foncteur $\varphi$ est \deleted{\unsure{conservatif}}\inserted{fidèle}]

\uline{Cor 1} Pour que $f:X\to Y$ dans $\mathcal M$ soit un mono, il f. et s. que $\varphi(f)$ le soit. L'application, pour $X\in\mathrm{Ob}\mathcal M$
\[
\mathrm{Ssobj}(X)\xrightarrow{\varphi}\mathrm{Ssobj}(\varphi(X))
\]
est injective, la relation d'ordre du premier ens est induite par \inserted{celle du second}

\uline{Déf} Sous-espace $(\mathcal M\text{-})$modéré de $|X|$ (stable par inters. finies)
Structure modérée \inserted{sur espace compact $K$,} Sous-espace modéré d'une structure modérée. Application modérée

\inserted{\uline{Cor 2} Une application \deleted{$\unsure{F}$} $f:|X|\longrightarrow|Y|$ modérée ssi son graphe l'est.}

M3) L'image d'un espace modéré par une application modérée est modérée

[
\diagram{\begin{tikzcd}[ampersand replacement=\&,row sep=small,column sep=small]
X\arrow[rr,"{f}",no head]\&\&Y\\
\&\times\&\\
X'\arrow[uu,hook]\arrow[rr]\&\&Y'\arrow[uu,hook]
\end{tikzcd}}
$(f:$
\diagram{\begin{tikzcd}[ampersand replacement=\&,row sep=small,column sep=small]
\lvert X\rvert\arrow[rr]\&\&\lvert Y\rvert\\
\&\times\&\\
\lvert X'\rvert\arrow[uu,hook]\arrow[rr]\&\&\lvert Y'\rvert\arrow[uu,hook]
\end{tikzcd}}
$)$

\margin{\uline{forme équivalente}}
M3\unsure{')} donc $f$ se factorise de façon \inserted{ess$^{\mathrm t}$} unique en
\[
X\xrightarrow{f'}Y'\overset{i}{-}Y
\]
avec $|f'|$ surjectif et $|i|$ injectif i.e. $i$ mono.

Est-ce que $f'$ épi (effectif) i.e.
\[
X\times_{Y'}X\rightrightarrows X-Y'\quad\text{exact ?}
\]
C'est épi car $|f'|$ l'est, et $|\;|$ conservatif,
soit alors $X\xrightarrow{g}Z$ tel que $g\,\mathrm{pr}_1=g\,\mathrm{pr}_2$, se factorise par $Y'$ ?
\deleted{$\unsure{X}$} $|g|$ se factorise en $h:|Y'|\to|Z|$, il suffit de voir le

\uline{Lemme} Soient $X\xrightarrow{f'}Y'$ dans $\mathcal M$ et $H:|Y'|-|Z|$ tel que $H\circ|f'|$ modérée, \inserted{i.e. $=|g|$,} $|f'|$ surjectif, alors $H$ modérée i.e. provient de $h:Y'\longrightarrow|Z|$.

En effet, \deleted{on voit d'abord} le graphe de $H$ est l'image du graphe de $|g|=H\circ f'$, donc est modéré.

\deleted{\unsure{M4} \illegible{1}}

\uline{Ind-objet spécial}\inserted{$*$} de $\mathcal M$ : a) définiss. \inserted{\unsure{par}} une \uline{suite}
\[
X_1-X_2\cdots-X_n-\cdots
\]
avec morphismes de transition des mono

b) $\forall i,\;\exists j\geq i$ tel que $\forall k\geq j,\quad |X_i|\subset\mathrm{Int}_{|X_k|}|X_j|$

\margin{55}
